\documentclass[12pt]{article}
\usepackage[utf8]{inputenc}
\usepackage[T1]{fontenc}
\usepackage{hyperref}
\title{Doutai Language Resources for K-12 Curriculum Development}
\author{Compiled by Claude Sonnet 5 (claude-sonnet-5) for Andr\'as Kornai}
\date{}
\begin{document}
\maketitle
\section{Introduction}
Doutai (ISO 639-3 tds; also Taori, Tolitai) is a Lakes Plain language of Indonesian Papua, spoken around Toli-Dou village near Taiyeve. Estimates conflict: about 70 speakers (Wikipedia, Endangered Languages citing Wurm 2007) versus 500 (Joshua Project). Endangered Languages reports no literacy in the language.
\section{Primary Educational Resources}
None found.
\section{Linguistic and Reference Materials}
\begin{itemize}
\item Wikipedia, Doutai language: \url{https://en.wikipedia.org/wiki/Doutai_language}
\item Endangered Languages Project entry: \url{https://www.endangeredlanguages.com/lang/1486}
\item Clouse 1997 on Lakes Plain classification (Trans New Guinea database): \url{https://transnewguinea.org/source/clouse-1997?page=87}
\end{itemize}
\section{Religious and Community Resources}
Joshua Project lists Bible translation as started, with no published materials.
\begin{itemize}
\item Joshua Project, Taori-So, Doutai in Indonesia: \url{https://www.joshuaproject.net/people_groups/15262/ID}
\end{itemize}
\section{Digital and Audio Resources}
None found.
\section{Curriculum Development Context}
The language is severely endangered and undocumented in teaching form. Documentation, orthography work and community decisions would precede any curriculum; Indonesian is the likely school language.
\end{document}